% ===== 原文 PDF pp.28--31（印刷 pp.27--30）=====
% 译者：tr-02；tag 区间 2.42--2.53；图 4、表 2、§3 标题与 §3 引子

% [原文 p.27]
近年来，对称性与 't Hooft 反常的概念被推广并重新表述，这是一个活跃的研究领域。在下文中，
我们将主要考虑连续 0-形式整体对称性（以及时空对称性）的微扰反常，并在结论中对 4d $\cN=2$
SCFT 里离散与广义整体对称性的非微扰反称作一些评论。

一般 $d$ 维 QFT 中 Lorentz 与整体对称性可能出现的微扰反常，由一个次数为 $d+2$ 的反常多项式
$\cI_{d+2}$ 分类，它是与这些对称性耦合的背景规范场的示性类的多项式。对 $d=4$，可能的反常是
混合引力-阿贝尔反常、纯非阿贝尔反常、混合阿贝尔-非阿贝尔反常与纯阿贝尔反常，分别对应下面
$\cI_6$ 中的四项：
\begin{equation}
  \cI_6=\alpha c_1(F)p_1(T)+\beta c_3(F)+\gamma c_2(F)c_1(F)+\delta c_1(F)^3\,,
  \tag{2.42}
\end{equation}
其中 $c_i(F)$ 记曲率为 $F$ 的背景规范场的第 $i$ 个陈类，$p_1(T)$ 是时空流形切丛的
Pontryagin 类。

任何 4d SCFT 都有一个由 $\mathrm{U}(1)_r$ R 对称性给出的特殊整体对称性，而相关的（混合）反常
由下面两个系数 $k_{rrr}$ 与 $k_r$ 支配，
\begin{equation}
  \cI_6\supset-\frac{k_r}{24}c_1(F)p_1(T)+\frac{k_{rrr}}{6}c_1(F)^3
  \tag{2.43}
\end{equation}
其中 $F$ 限制为上述 $\mathrm{U}(1)_r$ 背景。这些反常系数由涉及 $\mathrm{U}(1)_r$ 流与能动张量的
三点函数的宇称奇结构决定，如图 4 所示。\footnote{顺便说一句，一般的非共形 4d $\cN=2$
规范理论中，$\mathrm{U}(1)_r$ R 对称性由于理论中的动力学规范场而带有
Adler--Bell--Jackiw（ABJ）反常，其反常系数为 $2b$，其中 $b$ 是 (2.41) 所定义的一圈 beta
函数系数。因此，当 $b\neq0$ 时，该 R 对称性被明显破缺为 $\Z_{2b}$。反之，若这样的规范理论
是共形的，即 $b=0$，则 $\mathrm{U}(1)_r$ 没有 ABJ 反常，是一个名副其实的对称性——这正是
对一般 4d $\cN=2$ SCFT 所预期的。} 下面我们列出 4d $\cN=1$ 矢量多重态 $V$ 与手征多重态 $Q$
对 $(k_{rrr},k_r)$ 的贡献：
\begin{equation}
  \begin{aligned}
    V:\quad &k_{rrr}:=\Tr r^3_{\cN=1}=\rank G\,, & k_r&:=\Tr r_{\cN=1}=\rank G\,,\\
    Q:\quad &k_{rrr}:=\Tr r^3_{\cN=1}=(r(q)-1)^3\dim(Q)\,, & k_r&:=\Tr r_{\cN=1}=(r(q)-1)\dim(Q)\,,
  \end{aligned}
  \tag{2.44}
\end{equation}
其中迹 $\Tr$ 是对所有 Weyl 费米子求和的，$r(q)$ 记 4d $\cN=1$ 手征多重态 $Q$ 的标量分量的
$\mathrm{U}(1)_r$ R 荷。

% [Tutor Note: (2.42) 的四项与正文列出的四种反常按下标一一对应：$c_1(F)p_1(T)$ 是混合引力-阿贝尔（$p_1(T)$ 来自引力背景的切丛），$c_3(F)$ 是纯非阿贝尔，$c_2(F)c_1(F)$ 是混合阿贝尔-非阿贝尔，$c_1(F)^3$ 是纯阿贝尔。写 (2.43) 时之所以只剩两项，是因为此时 $F$ 限制为 $\mathrm{U}(1)_r$ 的背景：阿贝尔场没有 $c_2,c_3$，故只保留 $c_1(F)p_1(T)$ 与 $c_1(F)^3$。]
% [Tutor Note: (2.44) 中手征多重态 $Q$ 的贡献用 $(r(q)-1)^3$ 而不是 $r(q)^3$：$r(q)$ 是「标量」分量的 R 荷，而对所有「Weyl 费米子」求和时，同一 $\cN=1$ 多重态中费米子的 R 荷比标量少 1（R 对称性作用在超荷上，$\psi$ 与 $\phi$ 的 R 荷相差 1）。这一 $-1$ 是核对 (2.44) 与表 2 时最容易出错的地方。]

在 't Hooft 反常之外，偶数维的共形场论还带有 Weyl 反常：它们通过时空流形的曲率不变量修改
能动张量的无迹条件，因此也被称为迹反常。在 4d 中，迹反常取如下形式\footnote{一般偶数维 CFT
的迹反常取类似形式。特别地，$a$-反常对所有维数都是普适的，而 Weyl 部分则依赖于具体维数。
例如在 $d=2$ 时没有 Weyl 部分，而在 $d=6$ 时 Weyl 部分有三个独立项。}
\begin{equation}
  \langle T^{\mu}_{\ \mu}\rangle=\frac{c}{16\pi^2}(\mathrm{Weyl})^2-\frac{a}{16\pi^2}(\mathrm{Euler})\,,
  \tag{2.45}
\end{equation}
其中
\begin{equation}
  (\mathrm{Weyl})^2=R^2_{\mu\nu\rho\sigma}-2R^2_{\mu\nu}+\frac{1}{3}R^2\,,\qquad
  (\mathrm{Euler})=R^2_{\mu\nu\rho\sigma}-4R^2_{\mu\nu}+R^2\,.
  \tag{2.46}
\end{equation}

% [原文 p.28]
\begin{center}
  \includegraphics[width=0.8\linewidth]{fig4.png}
\end{center}
\figtitle{4}{由 $\mathrm{U}(1)_r$ 流 $J_r$ 与能动张量 $T$ 的三点函数给出的反常系数 $k_{rrr}$ 与 $k_r$。在拉格朗日理论中，它们由在 1 圈三角形中跑动的无质量 Weyl 费米子决定。}

常数系数 $(a,c)$ 称为 4d CFT 的共形中心荷，由 $\cN=1$ 超对称，它们可以用上述 $\mathrm{U}(1)_r$
反常表示 [138,139]：
\begin{equation}
  a=\frac{3}{32}\big[3k_{rrr}-k_r\big]\,,\qquad
  c=\frac{1}{32}\big[9k_{rrr}-5k_r\big]\,.
  \tag{2.47}
\end{equation}
之所以存在这样一个非平凡的关联，是因为 $\cN=1$ 超对称把 R 对称性流与能动张量联系起来。

% [Tutor Note: 图 4 内的英文标注是原图固有的：$J_r$ 表示 $\mathrm{U}(1)_r$ 流（虚线），$T$ 表示能动张量（点线），"odd" 指取三点函数的宇称奇部分。三个外线中恰有一个是 $T$（右图）时给 $k_r$，三个都是 $J_r$（左图）时给 $k_{rrr}$；这个「奇数个」结构正是 $R$ 反常的三角图特征。]
% [Tutor Note: (2.45) 左边的 $\langle T^{\mu}_{\ \mu}\rangle$ 是能动张量迹的期望值；$(\mathrm{Weyl})^2$ 与 $(\mathrm{Euler})$ 的具体组合由 (2.46) 给出，注意文中的 $R^2_{\mu\nu}$ 记号指 $R_{\mu\nu}R^{\mu\nu}$，$R^2_{\mu\nu\rho\sigma}$ 指 $R_{\mu\nu\rho\sigma}R^{\mu\nu\rho\sigma}$，$R^2$ 是 Ricci 标量的平方。(2.47) 则是把 (2.45) 的 $a,c$ 与 (2.43)(2.44) 的 't Hooft 系数连起来的关键一步，下一段对 $\cN=2$ 的 (2.48) 就是它的 $\cN=2$ 版本。]

我们这里关心的是 R 对称性为 $\mathrm{U}(1)_{r_{\cN=2}}\times\SU(2)_R$ 的 4d $\cN=2$ SCFT，
它们带有（混合）反常。如前，由 $\cN=2$ 超对称，这些反常可以用 Weyl 反常 $a$ 与 $c$ 表示 [75]，
\begin{equation}
  \resizebox{0.98\textwidth}{!}{$\displaystyle
  \cI_6\supset(c-a)c_1\big(F_{\mathrm{U}(1)_{r_{\cN=2}}}\big)p_1(T)+(a-c)c_1\big(F_{\mathrm{U}(1)_{r_{\cN=2}}}\big)\big)^3
  +2(2a-c)c_1\big(F_{\mathrm{U}(1)_{r_{\cN=2}}}\big)c_2\big(F_{\SU(2)_R}\big)\,.$}
  \tag{2.48}
\end{equation}
对上述结果的一个简单检验，是与 (2.43) 和 (2.47) 作比较，并注意 (2.18) 中 $\cN=1$ 与
$\cN=2$ 的 $\mathrm{U}(1)_r$ 对称性之间的关系。我们把 $\cN=2$ 自由场的反常汇总于表 2。

% [Tutor Note: (2.48) 中第二项在原文印作 $c_1(F_{\mathrm{U}(1)_{r_{\cN=2}}}))^3$，即在 $c_1(\cdot)$ 之后多出一个右括号（700 dpi 裁图复核 p.29 该式：同式前后两处 $c_1(F)$ 一处写 `))^3`、一处写 `)^3`，前后不对称）。译文照原文排印；正确形式应为 $c_1(F_{\mathrm{U}(1)_{r_{\cN=2}}})^3$，即与表 2 中矢量多重态那一格的第三项写法一致。]

% [原文 p.29]
\begin{center}
  \resizebox{0.99\linewidth}{!}{%
  \begin{tabular}{lccc}
    \toprule
    & $a$ & $c$ & \multicolumn{1}{c}{$\cI_6$}\\
    \midrule
    矢量多重态 & $\frac{5}{24}$ & $\frac{1}{6}$ &
    $-\frac{1}{24}c_1(F_{\mathrm{U}(1)_{r_{\cN=2}}})p_1(T)
     +\frac{1}{24}c_1(F_{\mathrm{U}(1)_{r_{\cN=2}}})^3
     +\frac{1}{2}c_1(F_{\mathrm{U}(1)_{r_{\cN=2}}})c_2(F_{\SU(2)_R})$\\
    \addlinespace
    超多重态 & $\frac{1}{24}$ & $\frac{1}{12}$ &
    $\frac{1}{24}c_1(F_{\mathrm{U}(1)_{r_{\cN=2}}})p_1(T)
     -\frac{1}{24}c_1(F_{\mathrm{U}(1)_{r_{\cN=2}}})^3$\\
    \bottomrule
  \end{tabular}}
\end{center}
\tabtitle{2}{4d $\cN=2$ 自由场的反常。}

% [Tutor Note: 表 2 的列归属逐格核对过：$a,c$ 两列是该自由场的共形中心荷，最右列是它代入 (2.48) 后对反常多项式 $\cI_6$ 的贡献（注意矢量多重态的第一项带负号 $-\frac{1}{24}$，超多重态的第一项带正号 $+\frac{1}{24}$，两者恰好相反）。用 (2.49) 可以自检这张表的读数：$n_v=1,n_h=0$ 给出 $a=5/24$、$c=2/12=1/6$，正是矢量多重态那一行；$n_v=0,n_h=1$ 给出 $a=1/24$、$c=1/12$，正是超多重态那一行。]

当 SCFT 具有共形拉格朗日量描述时（即 2.5 中 $b=0$ 的情形），'t Hooft 反常可直接从表 2 读出，
而共形中心荷简单地由下式给出
\begin{equation}
  a=\frac{5n_v+n_h}{24}\,,\qquad c=\frac{2n_v+n_h}{12}\,,
  \tag{2.49}
\end{equation}
其中 $n_v$ 与 $n_h$ 分别是理论中 $\cN=2$ 矢量多重态与超多重态的总数。

对非拉格朗日 SCFT，确定其反常一般要困难得多。所幸我们有一个强有力的工具——反常匹配
（anomaly matching），它利用反常在形变下的不变性。如果理论容许一个到弱耦合描述的对称形变，
反常就可以从那里被还原出来。在 $\cN=2$ SCFT 的语境下，有用的形变涉及走到理论的真空模空间上。
例如，在 SCFT 的 CB 上，$\mathrm{U}(1)_r$ 对称性与共形对称性被自发破缺，而 $\SU(2)_R$
得以保持。像通常那样，这意味着涉及 $\mathrm{U}(1)_r$ 与 $\SU(2)_R$ 的混合反常会同时收到
来自库仑分支 EFT 中场内容的贡献、以及来自含 Goldstone 玻色子的 Wess--Zumino 项的贡献。
利用 (2.48) 给出的 $(a,c)$ 与 't Hooft 反常之间的关系，就得到不动点处共形中心荷的如下公式 [75]，
\begin{equation}
  a=\frac{5r}{24}+\frac{h}{24}+\frac{R(A)}{4}+\frac{R(B)}{6}\,,\qquad
  c=\frac{r}{6}+\frac{h}{12}+\frac{R(B)}{3}\,,
  \tag{2.50}
\end{equation}
其中 $r$ 与 $h$ 分别计数该 4d SCFT 库仑分支一般点上的自由矢量多重态与自由超多重态的数目；
$R(A)$ 与 $R(B)$ 则捕捉来自 Wess--Zumino 项的贡献 [140--144]。$R(A)$ 与 $R(B)$ 可以从
耦合一般背景场的 CB EFT 得到，而它们可以用基本的 CB 数据——例如 CB 手征主算子的标度维数谱
——写出简洁的表达式 [75]（另见 [145] 的式 (1.1a)--(1.1c)）。我们将在 5.1 节中具体看到如何
为一大批 SCFT 提取 $R(A)$ 与 $R(B)$。

% [Tutor Note: (2.50) 与 (2.49) 的形式看起来相似，但地位不同：(2.49) 只在 $b=0$（有共形拉格朗日量、可微扰计算）时成立；而 (2.50) 是从 CB 上的反常匹配推出的，对非拉格朗日 SCFT 同样适用——代价是要额外知道 $R(A),R(B)$ 这两个 Wess--Zumino 项系数。原文强调 $R(A),R(B)$「能用 CB 手征主算子的标度维数谱表示」，这正是第 3 节（Seiberg--Witten 几何）与第 5 节（几何工程）能算出 SCFT 数据的接口。]
% [Tutor Note: 本段「反常匹配」的逻辑值得留意：形变要「对称」（symmetric deformation），即不破缺待匹配的整体对称性，否则 't Hooft 反常的匹配无从谈起；而在 CB 上 $\mathrm{U}(1)_r$ 与共形对称性自发破缺，所以必须把 Goldstone 玻色子的 Wess--Zumino 项一并计入，这正是 $R(A),R(B)$ 的来源。]

% [原文 p.30]
最后，如果该 $\cN=2$ SCFT 具有整体对称性 $G_{\mathrm{flavor}}$（为简单起见我们假定它是单群），
就可以再引入一个中心荷 $k_{\mathrm{flavor}}$，用来刻画混合的
$\mathrm{U}(1)_r$-$G_{\mathrm{flavor}}$ 反常，
\begin{equation}
  \cI_6\supset-\frac{1}{2}k_{\mathrm{flavor}}c_1(F_r)c_2(F_{\mathrm{flavor}})\,.
  \tag{2.51}
\end{equation}
与其他微扰反常一样，它由 $\mathrm{U}(1)_r$ 流与 $G_{\mathrm{flavor}}$ 流的三点函数中的宇称奇
结构决定。借助 $\cN=2$ 超对称，这个反常还与两个 $G_{\mathrm{flavor}}$ 流的 OPE 相关 [75,146]，
\begin{equation}
  \langle J^a_\mu(x)J^b_\nu(0)\rangle
  =\frac{3k_{\mathrm{flavor}}}{4\pi^4}\delta^{ab}\frac{x^2g_{\mu\nu}-2x_\mu x_\nu}{x^8}
  +\frac{2}{\pi^2}f^{abc}\frac{x_\mu x_\nu\,x\cdot J^c(0)}{x^6}+\dots
  \tag{2.52}
\end{equation}
注意这里流 $J^a_\mu$ 的归一化是由 RHS 第二项确定的，而结构常数 $f^{abc}$ 满足
\begin{equation}
  f^{aed}f^{bde}=2h^{\vee}\delta^{ab}\,,
  \tag{2.53}
\end{equation}
其中 $h^{\vee}$ 是 $G_{\mathrm{flavor}}$ 的对偶 Coxeter 数。在这一约定下，$n_H$ 个自由超多重态
具有 $G_{\mathrm{flavor}}=\mathrm{USp}(2n_H)$ 的味对称性，其反常系数 $k_{\mathrm{flavor}}=1$。
关于非平凡 SCFT 中味中心荷 $k_{\mathrm{flavor}}$ 的计算，可参见例如 [146,147]。

% [Tutor Note: (2.52) 的两个结构常数记号要对上：第一项（$\delta^{ab}$ 项）的系数给 $k_{\mathrm{flavor}}$，第二项（$f^{abc}$ 项）给的是流的归一化约定——原文紧接着说「$J^a_\mu$ 的归一化由 RHS 第二项确定」。因此 $k_{\mathrm{flavor}}=1$（自由超多重态）这个数值是「依赖该约定」的，比较不同文献的 $k_{\mathrm{flavor}}$ 前必须先对齐 (2.53) 的 $h^{\vee}$ 归一化（对 $\SU(N)$，$h^{\vee}=N$，$f^{aed}f^{bde}=N\delta^{ab}$ 即标准归一化）。]
% [Tutor Note: (2.53) 用对偶 Coxeter 数 $h^{\vee}$ 归一化结构常数，与 2.4 节脚注 15 中「$T(\mathrm{adj})$ 取为对偶 Coxeter 数」是同一套约定，二者应当互相一致：(2.37)(2.41) 里的 $T(\mathrm{adj})$ 与 (2.53) 里的 $h^{\vee}$ 对 $\SU(N)$ 都等于 $N$。]

\SEC{3}{库仑分支有效理论与 Argyres--Douglas 点}

在本节中，我们将讨论支配一般 4d $\cN=2$ 理论在 CB 模空间上低能动力学的库仑分支（CB）
有效场论（EFT）。我们将说明，这样的 EFT 为研究活在 CB 原点处的 4d $\cN=2$ SCFT 提供了
不可或缺的工具。为完整起见，先回顾 Seiberg--Witten（SW）的开创性工作 [20,21]、并对 SW 几何
作一介绍是有益的。正如我们将要说明的，4d $\cN=2$ 理论 CB EFT 的许多性质与动力学，都可以由
一个辅助的几何对象——所谓的 Seiberg--Witten 曲线——确定下来。

CB 上的基本自由度（其复维数为 $\dim_{\C}\cM_{\mathrm{CB}}=r$）是 $r$ 个 $\cN=2$ 阿贝尔
矢量多重态（其中包含参数化 $\cM_{\mathrm{CB}}$ 的复标量场）与由 Maxwell 作用支配的
$\mathrm{U}(1)^r$ 规范场，以及它们的费米伙伴。如果该理论在 UV 具有规范理论描述（如 2.3 节
所讨论），那么这些阿贝尔矢量多重态对应于秩为 $r$ 的规范群 $G$ 的 Cartan 生成元：沿真空模空间
的 CB 方向，超多重态标量取零值 $\langle q\rangle=\langle\tilde q\rangle=0$，而非阿贝尔矢量
多重态标量则获得 VEV $\langle\Phi\rangle$，这些 Cartan 生成元正是沿该分支被保留下来的那部分。

% [边界] 本句（PDF p.31 起于「If the theory has a gauge theory description in the UV ...
% then」，PDF p.32 续完）已写完，tr-03 从 PDF p.32 首句「For $G=\SU(r+1)$, up to a gauge
% transformation, ...」续。

% [Tutor Note: 本节开头把「CB EFT」与「CB 原点处的 SCFT」并置，这是全节的纲领：CB EFT 只在模空间的一般点上成立（那时只有 $r$ 个 $\mathrm{U}(1)$ 矢量多重态无质量，$\mathrm{U}(1)^r$ 有效描述成立）；在原点处低能谱会变大（出现额外的无质量态），CB EFT 本身不再有效，但它的 Seiberg--Witten 几何仍能编码原点 SCFT 的数据（3.4--3.5 节的做法）。读后面几节时务必区分「一般点上的 EFT」与「原点处的 SCFT」两种语境，$r$ 这个记号在前者指 CB 复维数、在 (2.50) 中指自由矢量多重态个数，恰是同一件事。]
